% Source: https://artofproblemsolving.com/community/c3h1596756_simplifying_a_somewhat_strange_expression
% Stable identity: aops-topic-1596756-q01
Multiply both sides by $(x-2)(x+3)$:
\[
5x=a(x+3)+b(x-2)=(a+b)x+(3a-2b)
\]
for all $x\ne2,-3$. Two polynomials that agree for infinitely many $x$ are identical, so we compare coefficients:
\[
a+b=5,\qquad 3a-2b=0 .
\]
From the second equation $b=\dfrac{3a}{2}$; then $a+\dfrac{3a}{2}=5$ gives $a=2$ and $b=3$. (Quicker: put $x=2$ to get $10=5a$, and $x=-3$ to get $-15=-5b$.)

So $\boxed{a=2,\ b=3}$:
\[
\dfrac{5x}{(x-2)(x+3)}=\dfrac{2}{x-2}+\dfrac{3}{x+3}.
\]
